% Research draft: method details and experimental claims must be updated
% after the proposed study has been implemented and evaluated.
\section{Introduction}
\label{sec:introduction}

Medical images vary across institutions because of differences in acquisition devices, imaging protocols, patient populations, and sample preparation~\cite{ali2023polypgen}. These variations create a persistent challenge for visual recognition: a model that performs well in its training environment may fail at an unseen medical center. Zech et al. document variable cross-hospital performance and hospital-associated confounding in pneumonia detection from chest radiographs~\cite{zech2018generalization}. The challenge extends beyond achieving high average accuracy. Reliable recognition requires representations that preserve diagnostically relevant visual structure while remaining robust to changes in how images are acquired. Multi-center datasets such as PolypGen provide an opportunity to investigate this problem under explicit institutional distribution shifts~\cite{ali2023polypgen}.

Pretrained visual encoders offer reusable representations for downstream recognition, including frozen-feature evaluation~\cite{oquab2024dinov2}. Their benefit in medical imaging is task-dependent: Raghu et al. found limited performance gains from ImageNet transfer in the two medical tasks they studied~\cite{raghu2019transfusion}. However, pretraining alone does not ensure that downstream decision rules generalize across centers. A classifier learned from pooled source data may exploit acquisition characteristics that correlate with class labels in the training set~\cite{zech2018generalization}. Theoretical results in domain adaptation also show that distribution alignment alone need not ensure low target error~\cite{zhao2019invariant}. Whether alignment suppresses useful variation in our medical setting remains an empirical question. This raises a central question: how can a recognition model retain class-discriminative structure without relying on center-specific visual cues?

Prototype-based learning provides a natural framework for studying this question. By representing classes through reference vectors in an embedding space, prototype methods make the relationship between representation geometry and prediction explicit~\cite{snell2017prototypical}. Strong cosine-classifier baselines further motivate controlled evaluation of whether a proposed prototype mechanism contributes beyond its underlying representation~\cite{chen2019closer}. Prior work represents categories through multiple prototypical image parts, as in ProtoPNet~\cite{chen2019protopnet}. Here, we study global embedding prototypes and hypothesize that a single prototype may compress distinct class appearances. The pooled mean also gives greater influence to centers contributing more samples, as formalized in Eq.~\ref{eq:pooled_centroid}. Introducing multiple prototypes increases representational flexibility, but does not automatically distinguish meaningful within-class variation from acquisition-specific structure. Consequently, neither pooled class centroids nor unconstrained collections of prototypes necessarily provide a decision rule that transfers to an unseen center.

We investigate \emph{domain-robust prototype learning for medical image recognition}. Our guiding hypothesis is that prototypes should capture variation that supports class discrimination while limiting dependence on the composition of the source centers. We study this hypothesis using pretrained image representations and controlled comparisons between pooled centroids, learned cosine classifiers, and multiple-prototype formulations. Center information is available during training to characterize source-domain variation, while inference must operate without access to the identity or labeled examples of the target center. This setting isolates generalization to unseen institutions from adaptation using target-domain data.

A rigorous evaluation protocol is essential to assess whether improvements reflect transferable visual structure. DomainBed emphasizes matched experimental conditions and source-only selection~\cite{gulrajani2021domainbed}; nested evaluation separates hyperparameter selection from performance estimation~\cite{varma2006bias}. We therefore formulate the study around held-out-center evaluation, with model selection restricted to source-domain validation data and related patient or video samples kept within the same partition. Comparisons use matched encoders and training conditions to separate the effect of prototype construction from the effect of backbone capacity or pretraining. Beyond aggregate performance, per-center results and controlled ablations examine when prototype learning improves generalization and when it removes useful class variation.

\paragraph{Scope and objectives.}
This paper focuses on supervised medical image classification under cross-center domain shift, using pretrained visual encoders and prototype-based decision rules. Models are trained on labeled source-center data and evaluated at centers excluded from training and model selection; neither target-center data nor target-center identity is used to adapt the model. PolypGen provides the initial experimental setting for polyp-presence classification. Within this scope, the study has three main objectives:
\begin{enumerate}
    \item \textbf{Develop a domain-robust prototype mechanism.} Investigate how center-balanced prototype estimation and multiple within-class prototypes can preserve class-discriminative variation while reducing sensitivity to source-center composition.
    \item \textbf{Evaluate generalization to unseen centers.} Compare the proposed mechanism with matched linear, cosine, nearest-centroid, and multiple-prototype baselines under held-out-center evaluation, reporting aggregate and worst-center performance with uncertainty estimates.
    \item \textbf{Explain the conditions governing prototype transfer.} Use controlled ablations and failure analysis to determine how prototype multiplicity, center balancing, and encoder adaptation affect generalization, including cases where balancing or prototype compression suppresses useful class variation.
\end{enumerate}
Together, these objectives connect prototype geometry to independently evaluated cross-center recognition and guide the design of lightweight medical image classifiers that remain effective beyond their training environments.

\section{Background and Related Work}
\label{sec:background}

\subsection{Cross-Center Domain Generalization}
Following the domain-generalization setting of DomainBed~\cite{gulrajani2021domainbed}, we model each imaging center as a domain with a joint distribution $P_s(X,Y)$ over images $X$ and labels $Y$. Given labeled datasets from $S$ source centers, the objective is to learn a predictor that transfers to an unseen target distribution $P_t(X,Y)$ sharing the same label space. Cross-hospital evaluation demonstrates differences in image appearance and class prevalence~\cite{zech2018generalization}, while class-conditional shift matters in analyses of transfer~\cite{zhao2019invariant}. We therefore treat center identity as a grouping variable rather than a complete description of acquisition style. In this paper, target-center images are unavailable during training and model selection. This is a domain-generalization setting~\cite{gulrajani2021domainbed}, whereas unsupervised domain adaptation permits access to unlabeled target images, as in the medical segmentation setting of En and Guo~\cite{en2024dpcl}.

Empirical risk minimization (ERM) on pooled source data is an essential baseline. Gulrajani and Lopez-Paz~\cite{gulrajani2021domainbed} demonstrate that careful implementation and model selection can make ERM competitive with dedicated domain-generalization algorithms on the benchmarks they examine. Their findings motivate matched architectures, tuning budgets, and source-only validation rather than assuming that an additional alignment objective is beneficial. Separately, Zhao et al.~\cite{zhao2019invariant} show, in domain adaptation, that distribution invariance and low source error do not suffice for successful transfer under certain shifts. These results caution against indiscriminate alignment; they do not prove that alignment necessarily harms the medical tasks considered here.

\subsection{Pretrained Representations and Prototype Geometry}
Let $f_\theta(x)$ be an image encoder and let $z(x)=f_\theta(x)/\|f_\theta(x)\|_2$ denote its normalized embedding, assuming a nonzero encoder output. Large-scale self-supervised encoders such as DINOv2 produce features transferable across multiple visual tasks~\cite{oquab2024dinov2}. We use such representations as a starting point rather than treating encoder pretraining as our contribution. A frozen-encoder experiment isolates the downstream decision rule, while a separately controlled adaptation experiment assesses whether changing the representation provides additional benefit.

Prototypical Networks classify examples by distances to class means in a learned metric space~\cite{snell2017prototypical}. For normalized class prototypes $p_c$ and temperature $\tau>0$, a cosine-based classifier instead uses
\begin{equation}
    q(y=c\mid x)=\frac{\exp(z(x)^\top p_c/\tau)}
    {\sum_{j=1}^{C}\exp(z(x)^\top p_j/\tau)}.
    \label{eq:prototype_classifier}
\end{equation}
With one trainable vector per class, this is a normalized linear classifier, consistent with the cosine formulation of Chen et al.~\cite{chen2019closer}. Supervised contrastive learning explicitly constructs labeled positive and negative relationships among embeddings~\cite{khosla2020supcon}. Cross-entropy training of cosine scores alone does not establish a new contrastive-learning method. The cosine baseline of Chen et al.~\cite{chen2019closer} reinforces the importance of testing prototype updates against a strong classifier using the same encoder.

ProtoPNet represents categories through prototypical image parts and exposes supporting regions~\cite{chen2019protopnet}. Its spatial, example-based explanation differs from our global embedding vectors. Spherical clustering provides an established mechanism for representing several directions in normalized feature space~\cite{dhillon2001spherical}. Whether these directions capture transferable class structure or acquisition artifacts is a hypothesis to test. A geometric reference vector alone is not a clinically validated explanation; interpretability requires evidence linking matches to meaningful image content. Temperature scaling is likewise an established calibration technique~\cite{guo2017calibration}, not a guarantee of calibration after domain shift.

\subsection{Prototype Weighting and Medical Transfer}
For class $c$, let $n_{s,c}$ and $\mu_{s,c}$ denote the sample count and mean embedding at source center $s$. Before normalization, the pooled class centroid is
\begin{equation}
    \mu_c^{\mathrm{pool}}=
    \frac{\sum_{s=1}^{S}n_{s,c}\mu_{s,c}}
    {\sum_{s=1}^{S}n_{s,c}},
    \label{eq:pooled_centroid}
\end{equation}
where centers with no examples of class $c$ contribute zero. Thus, centers with larger class-specific sample counts receive greater weight. Averaging the available center means equally changes the weighting, but is not necessarily preferable: estimates from small centers may be noisy, and equal weighting may differ from the deployment distribution. The effects of center balancing and prototype multiplicity are consequently empirical questions, not assumed improvements. Equation~\ref{eq:pooled_centroid} is an algebraic decomposition of the pooled mean, rather than a new estimator; our clustering also builds on established spherical $k$-means~\cite{dhillon2001spherical}.

Prototype-based medical transfer has existing precedents. En and Guo~\cite{en2024dpcl} introduce dynamic cross-domain prototypes and contrastive learning for medical image segmentation using labeled source and unlabeled target data. Their method addresses adaptation with target access, while our proposed study concerns classification without target access. This distinction defines a different evaluation setting; it does not establish novelty for prototype alignment itself. PolypGen supplies real multi-center endoscopy data and an existing out-of-sample assessment protocol~\cite{ali2023polypgen}. Because it includes video frames, partitions must respect available patient or sequence identifiers to reduce dependence between training and evaluation samples.

\paragraph{Research question and limits.}
The question pursued here is whether center-aware prototype estimation preserves useful within-class structure and improves unseen-center recognition beyond matched ERM, cosine, and centroid baselines. The reviewed work motivates this investigation but does not establish that the exact proposed mechanism is novel; a broader comparison with recent prototype-based domain-generalization methods remains necessary. Evidence must distinguish the effects of weighting, prototype count, and encoder adaptation, and include cases where the method fails. Results on PolypGen alone would support conclusions about the evaluated endoscopy setting, rather than universal robustness across medical modalities or readiness for clinical deployment.

\section{Methodology}
\label{sec:method}
% Proposed implementation specification, not a description of completed runs.
% Replace this comment and reconcile all settings with the released code
% before submission. No performance or novelty claim is established here.

\subsection{Setting and Design Rationale}
We formulate the proposed method as \emph{center-balanced multi-prototype classification} (CBMP). The training set contains tuples $(x_i,y_i,s_i,g_i)$, where $s_i$ denotes the source center and $g_i$ identifies a patient or video group. The target center is excluded from feature fitting, prototype construction, and hyperparameter selection. All centers share a fixed label vocabulary; every evaluated class must occur in the source training set. For PolypGen classification, labels indicate polyp presence, subject to an audited image-label manifest rather than assumptions based solely on folder names.

The core experiment freezes a pretrained encoder $f_\theta$ and learns only class prototypes. This choice separates prototype construction from representation adaptation. CBMP changes the influence of centers in class-conditional clustering and represents each class with several reference directions. It imposes no explicit domain-invariance loss and makes no assumption that center identity fully explains visual variation. Whether its weighting improves transfer is tested against identically encoded baselines.

Figure~\ref{fig:cbmp_system} summarizes the system and its experimental data boundaries. The same outer center split and encoder features are used for CBMP and its controls. Inner source-center validation selects the configuration before the outer source data are used to construct the final prototype bank. Target images enter only the fixed inference path.

% Add to the CVPR preamble: \usepackage{tikz}
% \usetikzlibrary{arrows.meta,positioning,shapes.geometric,shapes.symbols,fit,backgrounds}
\begin{figure*}[t]
\centering
\begin{tikzpicture}[
    font=\footnotesize, >=Latex,
    block/.style={draw=blue!65!black, fill=blue!6, rounded corners=3pt,
        align=center, text width=3.8cm, minimum height=1.35cm, inner sep=5pt},
    select/.style={block, draw=orange!75!black, fill=orange!12},
    learned/.style={block, draw=teal!65!black, fill=teal!10},
    target/.style={block, draw=purple!65!black, fill=purple!7, dashed},
    data/.style={block, shape=cylinder, shape border rotate=90,
        aspect=0.5, rounded corners=0pt},
    flow/.style={->, thick, draw=blue!65!black},
    choice/.style={->, thick, dashed, draw=orange!80!black},
    testflow/.style={->, thick, draw=purple!65!black}
]
\node[data] (source) at (0,0) {\textbf{Source centers}\\Labeled images\\Center and group IDs};
\node[select] (validation) at (5.2,0) {\textbf{Source-only validation}\\Select $K$, $\beta$, and $\tau$};
\node[data, draw=purple!65!black, fill=purple!7, dashed] (targetdata) at (10.4,0) {\textbf{Held-out target center}\\Excluded from fitting\\and model selection};

\node[block] (encoder) at (0,-2.6) {\textbf{Frozen encoder}\\Prescribed preprocessing\\Normalized embeddings $z_i$};
\node[select] (weight) at (5.2,-2.6) {\textbf{Group/center weights}\\$w_i\propto G_{s_i,c}^{-\beta}/m_{g_i,c}$\\Source training data only};
\node[target] (targetencoder) at (10.4,-2.6) {\textbf{Target encoding}\\Same frozen encoder\\No target adaptation};

\node[learned] (cluster) at (0,-5.2) {\textbf{Spherical clustering}\\Weighted class-wise updates\\Training-objective restarts};
\node[learned, double, double distance=1pt] (bank) at (5.2,-5.2) {\textbf{Stored prototype bank}\\$\{p_{c,k}\}$ and fixed settings\\Built from outer sources};
\node[target] (prediction) at (10.4,-5.2) {\textbf{Prediction}\\$h_c(x)=\max_k z(x)^\top p_{c,k}$\\Source-selected temperature};

\node[block, draw=gray!70, fill=gray!8, text width=8.5cm]
    (controls) at (2.6,-7.8) {\textbf{Matched controls and ablations}\\Same outer splits and encoder\\Isolate prototype count and weighting};
\node[target] (evaluation) at (10.4,-7.8) {\textbf{Independent evaluation}\\Per-center / worst-center metrics\\Group-level uncertainty};

\draw[flow] (source) -- (encoder);
\draw[choice] (source) -- (validation);
\draw[flow] (encoder) -- (weight);
\draw[choice] (validation) -- (weight);
\draw[flow] (weight.south) -- (5.2,-3.9) -- (0,-3.9) -- (cluster.north);
\draw[flow] (cluster) -- (bank);
\draw[flow] (bank) -- (prediction);
\draw[testflow] (targetdata) -- (targetencoder);
\draw[testflow] (targetencoder) -- (prediction);
\draw[testflow] (prediction) -- (evaluation);
\draw[flow] (cluster.south) -- (controls.north -| cluster.south);
\draw[flow] (controls) -- (evaluation);
\node[font=\scriptsize, align=center, text width=14cm] at (5.2,-9.35)
    {\textcolor{blue!65!black}{Blue: data/encoding}\quad
     \textcolor{orange!80!black}{Orange: source-only selection/weighting}\\[3pt]
     \textcolor{teal!65!black}{Teal: prototype construction}\quad
     \textcolor{purple!65!black}{Purple dashed: target evaluation}};
\end{tikzpicture}
\caption{CBMP system and evaluation boundaries. Cylinders denote datasets, rounded rectangles denote processing, and the double border identifies stored prototypes. Source-only inner validation selects the configuration, after which final prototypes are constructed from the outer sources. The dashed target path uses fixed encoding and prediction, with no feedback from target evaluation into fitting or selection. Controls share the outer splits and encoder; ablations isolate multiplicity and weighting. Colors reinforce the stage labels and boundary styles, so the diagram remains interpretable in grayscale.}
\label{fig:cbmp_system}
\end{figure*}

\subsection{Embeddings and Sampling Units}
We extract normalized embeddings $z_i=z(x_i)$ using the checkpoint's prescribed image preprocessing, in evaluation mode and without stochastic augmentation. Numerically zero embeddings are flagged as invalid rather than normalized. Preprocessing, checkpoint version, and feature extraction settings are fixed across methods. Any fitted transformation is estimated on source training data only.

Repeated frames should not determine a center's influence merely through sequence length. For class $c$, let $\mathcal{G}_{s,c}$ be the groups containing that class at center $s$, let $G_{s,c}=|\mathcal{G}_{s,c}|$, and let $m_{g,c}$ count class-$c$ training images in group $g$. For $\mathcal{I}_c=\{i:y_i=c\}$, define
\begin{equation}
    a_i=\frac{G_{s_i,c}^{-\beta}}{m_{g_i,c}},
    \qquad
    w_i=\frac{a_i}{\sum_{j\in\mathcal{I}_c}a_j},
    \label{eq:cbmp_weights}
\end{equation}
where $\beta\in[0,1]$ controls center balancing. Each group receives equal total weight within its center and class. A center's total unnormalized class weight is $G_{s,c}^{1-\beta}$: $\beta=0$ retains group-count weighting, whereas $\beta=1$ equalizes the available centers. Absent center-class combinations contribute no samples. Patient groups take precedence over sequence groups when available; the grouping rule and missing identifiers are documented. This weighting reduces frame-count dominance but does not make correlated observations statistically independent.

\subsection{Weighted Spherical Prototype Construction}
Building on spherical $k$-means~\cite{dhillon2001spherical}, for each class $c$ we construct $K_c$ unit-norm prototypes $p_{c,k}$ by minimizing the weighted spherical clustering objective
\begin{equation}
    J_c=\sum_{i\in\mathcal{I}_c}w_i
    \min_{1\leq k\leq K_c}\bigl(1-z_i^\top p_{c,k}\bigr),
    \qquad \|p_{c,k}\|_2=1.
    \label{eq:cbmp_clustering}
\end{equation}
Given prototypes, each embedding is assigned to the largest cosine similarity, with ties resolved by prototype index. For assigned set $\mathcal{A}_{c,k}$, the update is
\begin{equation}
    v_{c,k}=\sum_{i\in\mathcal{A}_{c,k}}w_i z_i,
    \qquad p_{c,k}=\frac{v_{c,k}}{\|v_{c,k}\|_2}.
    \label{eq:cbmp_update}
\end{equation}
This update minimizes the fixed-assignment objective whenever $v_{c,k}$ is nonzero. It does not guarantee a globally optimal clustering or target-domain robustness.

Initialization selects distinct training embeddings without replacement using probabilities proportional to $w_i$. We use five restarts per class and retain the solution with the smallest training objective. Empty clusters or numerically zero update vectors are reinitialized from the training embedding with the largest weighted residual to its nearest valid prototype, avoiding duplicate seeds where possible. Iteration stops when assignments are unchanged or after 100 updates. We choose $K_c=\min(K,|\mathcal{I}_c|,\sum_s G_{s,c})$ and record the realized prototype count. A class with fewer than $K_c$ distinct embedding directions uses fewer prototypes. These rules prevent unsupported or undefined prototype updates; they do not substitute for adequate class coverage.

\subsection{Prediction and Calibration}
For an unseen image, the class score is the strongest match among its prototypes:
\begin{equation}
    h_c(x)=\max_{1\leq k\leq K_c}z(x)^\top p_{c,k},
    \qquad \hat y(x)=\mathop{\rm argmax}_{c} h_c(x).
    \label{eq:cbmp_prediction}
\end{equation}
This decision rule lets different prototypes support different appearances without averaging them into a single direction. Prototype count can also increase accidental matches, motivating a matched multiple-prototype baseline. When $K=1$ and $\beta=0$, the method reduces to a group-weighted normalized class centroid; when $\beta=1$, it equalizes available centers before normalization.

For probabilistic evaluation, scores are converted to a softmax with temperature $\tau>0$, as in Eq.~\ref{eq:prototype_classifier}, substituting $h_c(x)$ for the single-prototype similarity. Following temperature scaling~\cite{guo2017calibration}, temperature is selected using source validation negative log-likelihood; it does not change the argmax decision. These scores are not assumed to remain calibrated after a center shift. Any operating threshold is also fixed from source validation data. Inference uses only the image, encoder, and prototypes, with no target-center identifier or batch adaptation.

\subsection{Source-Only Selection and Controlled Evaluation}
For each outer held-out center, we select $(K,\beta)$ through inner leave-one-source-center-out validation, using the candidate grid $K\in\{1,2,4,8\}$ and $\beta\in\{0,0.5,1\}$. The selection score is mean balanced accuracy across the inner validation centers. Ties favor smaller $K$, then smaller $\beta$. Temperature is chosen from $\{0.02,0.05,0.1,0.2,0.5,1\}$ using the mean inner-validation negative log-likelihood for the selected configuration. Inner folds must retain all classes in training; unsupported folds are reported and excluded by a rule fixed before evaluation. If no valid center-level inner fold exists, a documented group-disjoint source validation split is used. Once settings are fixed, prototypes are reconstructed from all outer source training data. Outer target scores never select settings, restarts, or checkpoints, consistent with the model-selection concerns raised by DomainBed~\cite{gulrajani2021domainbed}.

Controls include matched linear and cosine classifiers, a single-centroid classifier, and unbalanced multiple-prototype clustering. A factorial comparison of $K=1$ versus $K>1$ and $\beta=0$ versus $\beta>0$ separates multiplicity from center weighting. Group weighting is held constant in these comparisons; a separate image-weighted control measures its effect. Trainable classifiers receive the same source validation protocol and a disclosed tuning budget. Encoder adaptation, if evaluated, is a separate experiment with matched optimization and feature recomputation, rather than an unreported change to the frozen-encoder method.

We report balanced accuracy, macro F1, AUROC where defined, and per-center performance, with the worst-center score computed across the outer evaluation folds. Group-level bootstrap intervals resample the highest available independent unit; repeated optimization seeds are reported separately. Evaluation manifests, class coverage, seeds, selected settings, and predictions are saved so each result can be reconstructed. Center balancing cannot compensate for missing disease appearances, erroneous labels, or an encoder insensitive to the relevant visual cues.

\paragraph{Computational cost.}
With cached $d$-dimensional embeddings, $R$ restarts and $T$ clustering updates require approximately $O(RTNKd)$ operations for $N$ training images, excluding encoder feature extraction. Prototype storage is $O(d\sum_c K_c)$, and prediction costs $O(d\sum_c K_c)$ per image after encoding. The frozen-encoder construction uses no gradient updates. The proposed contribution must therefore be assessed through its weighting rule and independently measured transfer behavior, rather than architectural scale or a renamed cosine classifier.
